\documentclass[11pt]{article}
\usepackage{mathblatt}
\begin{document}
\weit
\typeout{PROBE-BEGIN symmetrie-abbildungen-zone-f2-v1}
\begin{aufgabe}{\detokenize{symmetrie-abbildungen-zone-f2-v1}}
\begin{teile}
\teil Probe

\geradenkreuzung{90}{\alpha}{\beta}{\gamma}{\delta}
\end{teile}
\end{aufgabe}
\typeout{PROBE-END symmetrie-abbildungen-zone-f2-v1}
\typeout{PROBE-BEGIN winkel-dreiecke-e2-k2-s0-v1}
\begin{aufgabe}{\detokenize{winkel-dreiecke-e2-k2-s0-v1}}
\begin{teile}
\teil Probe

\geradenkreuzung{55}{\alpha}{\beta}{\gamma}{\delta}
\end{teile}
\end{aufgabe}
\typeout{PROBE-END winkel-dreiecke-e2-k2-s0-v1}
\typeout{PROBE-BEGIN winkel-dreiecke-e2-k2-s0-v2}
\begin{aufgabe}{\detokenize{winkel-dreiecke-e2-k2-s0-v2}}
\begin{teile}
\teil Probe

\geradenkreuzung{64}{\alpha}{\beta}{\gamma}{\delta}
\end{teile}
\end{aufgabe}
\typeout{PROBE-END winkel-dreiecke-e2-k2-s0-v2}
\typeout{PROBE-BEGIN winkel-dreiecke-e2-k2-s1-v1}
\begin{aufgabe}{\detokenize{winkel-dreiecke-e2-k2-s1-v1}}
\begin{teile}
\teil Probe

\geradenkreuzung{38}{\alpha}{}{}{}
\end{teile}
\end{aufgabe}
\typeout{PROBE-END winkel-dreiecke-e2-k2-s1-v1}
\typeout{PROBE-BEGIN winkel-dreiecke-e2-k2-s1-v2}
\begin{aufgabe}{\detokenize{winkel-dreiecke-e2-k2-s1-v2}}
\begin{teile}
\teil Probe

\geradenkreuzung{51}{\alpha}{}{}{}
\end{teile}
\end{aufgabe}
\typeout{PROBE-END winkel-dreiecke-e2-k2-s1-v2}
\end{document}
